% =====================================================================
\chapter{k-Nearest Neighbours}
% =====================================================================

\labinfo{01\_KNN}{knn.py}{UCI Wine: 178 wines, 13 chemical measurements, 3 cultivars}
{Multi-class classification}{Test accuracy 92.5\%, macro F1 0.924 (k = 17, Manhattan distance)}

\section{Dataset}
The \textbf{Wine} dataset (UCI) contains the results of a chemical analysis of 178 wines grown in the
same region of Italy but from three different cultivars. Each wine is described by 13 continuous
measurements: alcohol, malic acid, ash, alcalinity of ash, magnesium, total phenols, flavanoids,
non-flavanoid phenols, proanthocyanins, colour intensity, hue, OD280/OD315 of diluted wines and
proline. The classes are a little unbalanced: 59, 71 and 48 wines.

The features live on very different scales (magnesium is around 70--162, proline 278--1680, while hue
and non-flavanoid phenols are below 2), which makes this dataset a good test of why k-NN needs scaling.
I split it with a stratified 70/30 split into 125 training and 53 test wines.

\section{Code and experiments}
The script \code{knn.py} is organised as follows.
\begin{itemize}
  \item \code{stratified\_split} and \code{stratified\_folds} split indices while keeping the class mix.
  \item \code{Standardizer} learns the mean and standard deviation on the training rows and applies them
        to any other rows.
  \item \code{pairwise\_distance} computes all query-to-training distances at once. For Euclidean distance
        it uses the identity $\|\vect a-\vect b\|^2=\|\vect a\|^2+\|\vect b\|^2-2\,\vect a^\top\vect b$, which turns the
        whole job into one matrix product.
  \item \code{KNNClassifier} has the usual \code{fit}/\code{predict} pair. \code{fit} only stores the data.
        \code{predict} sorts each row of the distance matrix, takes the first $k$ indices and runs the
        vote, with the nearest-neighbour tie-break described above.
  \item \code{confusion\_matrix} and \code{scores\_from\_confusion} compute accuracy and per-class
        precision, recall and F1.
\end{itemize}

\noindent The experiment has four parts:
(1)~5-fold cross-validation on the training set for $k=1,\dots,25$ with three settings: Euclidean with
scaling, Manhattan with scaling, and Euclidean without scaling;
(2)~the best $(k,\text{metric})$ pair is trained on the full training set and scored once on the test
set;
(3)~the same model is tried without scaling and with distance-weighted votes;
(4)~to see what the decision boundary looks like, I project the data onto its first two principal
components and colour the plane by the k-NN prediction.

\section{Results}

The cross-validation curves say the most important thing straight away: \textbf{scaling is not optional}.
Without it the CV accuracy stays between 71\% and 77\% whatever $k$ is; with it, accuracy jumps to
about 95--99\%. With scaled features both distances work well; Manhattan is slightly more stable for
small $k$. The best CV score was 99.2\% for Manhattan distance with $k=17$, so that is the final model.

\begin{table}[H]
\centering\small
\caption{Test-set results of the final model (Manhattan distance, $k=17$, 53 test wines).}
\label{tab:knn-test}
\begin{tabular}{@{}l S[table-format=1.3] S[table-format=1.3] S[table-format=1.3] S[table-format=2.0]@{}}
\toprule
\textbf{Class} & {\textbf{Precision}} & {\textbf{Recall}} & {\textbf{F1}} & {\textbf{Support}}\\
\midrule
Cultivar 1 & 0.947 & 1.000 & 0.973 & 18 \\
Cultivar 2 & 1.000 & 0.810 & 0.895 & 21 \\
Cultivar 3 & 0.824 & 1.000 & 0.903 & 14 \\
\midrule
Macro average & 0.924 & 0.937 & 0.924 & 53 \\
\midrule
\multicolumn{5}{@{}l}{Accuracy: \textbf{0.925} (49 of 53 correct)}\\
\multicolumn{5}{@{}l}{Same $k$ without scaling: 0.679 \quad Distance-weighted voting: 0.925}\\
\bottomrule
\end{tabular}
\end{table}

The test accuracy (92.5\%) is lower than the CV estimate (99.2\%). With only 53 test wines, every single
mistake costs almost 2 percentage points, so 4 mistakes instead of 0 or 1 is enough to explain the gap;
it is a reminder that small test sets give noisy numbers. Distance weighting did not change any
prediction here. Dropping the scaling, on the other hand, cost 25 percentage points on the test set as
well.

\begin{figure}[H]
\centering
\includegraphics[width=0.48\linewidth]{\figpath{01\_KNN}{cv_accuracy_vs_k.pdf}}\hfill
\includegraphics[width=0.48\linewidth]{\figpath{01\_KNN}{confusion_matrix.pdf}}
\caption{k-NN cross-validation curve and test confusion matrix.}
\end{figure}

\section{Source code}
\lstinputlisting[caption={\texttt{01\_KNN/knn.py}}]{01_KNN/knn.py}
\clearpage
